\documentclass[12pt,a4paper]{article}
\usepackage[UTF8]{ctex}
\usepackage{amsmath,amssymb,amsthm}
\usepackage{geometry}
\usepackage{graphicx}

\geometry{margin=2.5cm}

\title{导纳控制中$x$和$\dot{x}$的含义：当前状态 vs. 误差项}
\author{第11章补充说明}
\date{}

\begin{document}

\maketitle

\section{关键问题}

在导纳控制中，公式为：
\begin{equation}
M\ddot{x}_d + B\dot{x} + Kx = f_{\text{ext}} \tag{11.64}
\end{equation}

\textbf{重要问题}：公式中的$x$和$\dot{x}$是什么？

\textbf{答案}：它们是\textbf{当前实际状态}，不是误差项！

\section{导纳控制公式的详细解释}

\subsection{公式(11.64)的含义}

\begin{equation}
M\ddot{x}_d + B\dot{x} + Kx = f_{\text{ext}}
\end{equation}

\textbf{各项的含义}：
\begin{itemize}
\item $x$: \textbf{当前实际}末端执行器位置（由编码器测量）
\item $\dot{x}$: \textbf{当前实际}末端执行器速度（由编码器或转速计测量）
\item $\ddot{x}_d$: \textbf{期望的}末端执行器加速度（这是要计算的）
\item $M, B, K$: 虚拟质量、阻尼和刚度矩阵
\item $f_{\text{ext}}$: 感测到的外力（由力-力矩传感器测量）
\end{itemize}

\subsection{为什么使用当前状态？}

导纳控制的核心思想是：
\begin{itemize}
\item \textbf{感测外力}：测量用户施加的力$f_{\text{ext}}$
\item \textbf{根据当前状态计算响应}：基于当前位置$x$和速度$\dot{x}$，计算期望加速度$\ddot{x}_d$
\item \textbf{实现阻抗行为}：使机器人表现得像具有特定阻抗的系统
\end{itemize}

\textbf{物理意义}：
\begin{itemize}
\item $Kx$: 虚拟弹簧力，与\textbf{当前位置}成正比
\item $B\dot{x}$: 虚拟阻尼力，与\textbf{当前速度}成正比
\item $M\ddot{x}_d$: 虚拟惯性力，与\textbf{期望加速度}成正比
\item 这些力的总和等于外力$f_{\text{ext}}$
\end{itemize}

\section{与阻抗控制的区别}

\subsection{阻抗控制}

在阻抗控制中，控制律可能是：
\begin{equation}
\tau = J^T(\theta)\left(\tilde{\Lambda}(\theta)\ddot{x} + \tilde{\eta}(\theta, \dot{x}) - (M\ddot{x} + B\dot{x} + Kx)\right)
\end{equation}

这里：
\begin{itemize}
\item $x, \dot{x}, \ddot{x}$: 也是当前实际状态
\item 但控制目标是产生力，使系统呈现阻抗特性
\end{itemize}

\subsection{导纳控制}

在导纳控制中：
\begin{equation}
M\ddot{x}_d + B\dot{x} + Kx = f_{\text{ext}}
\end{equation}

这里：
\begin{itemize}
\item $x, \dot{x}$: \textbf{当前实际状态}（测量值）
\item $\ddot{x}_d$: \textbf{期望加速度}（计算值）
\item 控制目标是产生运动，使系统呈现导纳特性
\end{itemize}

\section{为什么不是误差项？}

\subsection{常见误解}

有人可能认为公式应该是：
\begin{equation}
M\ddot{x}_d + B(\dot{x}_d - \dot{x}) + K(x_d - x) = f_{\text{ext}} \quad \text{（错误！）}
\end{equation}

\textbf{这是错误的！}原因如下：

\subsection{正确的理解}

导纳控制的目标是：\textbf{使机器人表现得像具有特定阻抗的系统}。

\textbf{阻抗系统的特性}：
\begin{itemize}
\item 弹簧力：$F_{\text{spring}} = Kx$（与当前位置成正比）
\item 阻尼力：$F_{\text{damp}} = B\dot{x}$（与当前速度成正比）
\item 惯性力：$F_{\text{inertia}} = M\ddot{x}$（与当前加速度成正比）
\end{itemize}

\textbf{注意}：这些力都是基于\textbf{当前状态}，不是基于误差！

\subsection{物理类比}

想象一个质量-弹簧-阻尼器系统：
\begin{equation}
m\ddot{x} + b\dot{x} + kx = f
\end{equation}

这里：
\begin{itemize}
\item $x$: 质量的\textbf{当前位置}
\item $\dot{x}$: 质量的\textbf{当前速度}
\item $\ddot{x}$: 质量的\textbf{当前加速度}
\item 这些都是实际状态，不是误差
\end{itemize}

导纳控制就是要让机器人表现得像这样的系统，所以使用当前状态是合理的。

\section{具体例子}

\subsection{例子1：用户推机器人}

\textbf{场景}：
\begin{itemize}
\item 机器人当前位置：$x = 0.1$ m
\item 机器人当前速度：$\dot{x} = 0$ m/s
\item 用户施加力：$f_{\text{ext}} = 10$ N
\item 虚拟参数：$M = 1$ kg，$B = 2$ N·s/m，$K = 100$ N/m
\end{itemize}

\textbf{计算期望加速度}：
\begin{align}
M\ddot{x}_d + B\dot{x} + Kx &= f_{\text{ext}} \\
1 \cdot \ddot{x}_d + 2 \cdot 0 + 100 \cdot 0.1 &= 10 \\
\ddot{x}_d + 10 &= 10 \\
\ddot{x}_d &= 0 \text{ m/s}^2
\end{align}

\textbf{解释}：
\begin{itemize}
\item 虚拟弹簧力$Kx = 100 \times 0.1 = 10$ N（与当前位置成正比）
\item 虚拟阻尼力$B\dot{x} = 2 \times 0 = 0$ N（与当前速度成正比）
\item 外力$f_{\text{ext}} = 10$ N
\item 弹簧力正好抵消外力，所以期望加速度为0
\end{itemize}

\textbf{注意}：这里使用的是\textbf{当前实际位置}$x = 0.1$ m，不是误差！

\subsection{例子2：机器人运动中的响应}

\textbf{场景}：
\begin{itemize}
\item 机器人当前位置：$x = 0.2$ m
\item 机器人当前速度：$\dot{x} = 0.5$ m/s（向右运动）
\item 用户施加力：$f_{\text{ext}} = 5$ N（向右）
\item 虚拟参数：$M = 1$ kg，$B = 2$ N·s/m，$K = 100$ N/m
\end{itemize}

\textbf{计算期望加速度}：
\begin{align}
M\ddot{x}_d + B\dot{x} + Kx &= f_{\text{ext}} \\
1 \cdot \ddot{x}_d + 2 \cdot 0.5 + 100 \cdot 0.2 &= 5 \\
\ddot{x}_d + 1 + 20 &= 5 \\
\ddot{x}_d &= -16 \text{ m/s}^2
\end{align}

\textbf{解释}：
\begin{itemize}
\item 虚拟弹簧力$Kx = 100 \times 0.2 = 20$ N（向左，因为$x > 0$）
\item 虚拟阻尼力$B\dot{x} = 2 \times 0.5 = 1$ N（向左，因为$\dot{x} > 0$）
\item 外力$f_{\text{ext}} = 5$ N（向右）
\item 净力：$5 - 20 - 1 = -16$ N（向左）
\item 所以期望加速度为$-16$ m/s$^2$（向左减速）
\end{itemize}

\textbf{注意}：这里使用的是\textbf{当前实际状态}$x = 0.2$ m和$\dot{x} = 0.5$ m/s！

\section{如果使用误差项会怎样？}

\subsection{错误的公式}

如果错误地使用误差项：
\begin{equation}
M\ddot{x}_d + B(\dot{x}_d - \dot{x}) + K(x_d - x) = f_{\text{ext}} \quad \text{（错误！）}
\end{equation}

\textbf{问题}：
\begin{enumerate}
\item \textbf{缺少参考点}：$x_d$是什么？在导纳控制中，通常没有明确的期望位置$x_d$
\item \textbf{物理意义不清}：这不再是阻抗系统的行为
\item \textbf{控制目标不同}：这变成了轨迹跟踪，而不是阻抗模拟
\end{enumerate}

\subsection{正确的理解}

导纳控制的目标是：
\begin{itemize}
\item \textbf{不是}跟踪特定轨迹（所以不需要$x_d$）
\item \textbf{而是}响应外力，表现得像阻抗系统
\item 因此使用当前状态是合理的
\end{itemize}

\section{与轨迹跟踪控制的对比}

\subsection{轨迹跟踪控制}

在轨迹跟踪控制中，我们使用误差项：
\begin{equation}
\dot{\theta} = \dot{\theta}_d + K_p(\theta_d - \theta)
\end{equation}

这里：
\begin{itemize}
\item $\theta_d$: 期望位置（由轨迹规划给出）
\item $\theta$: 实际位置
\item $\theta_d - \theta$: 位置误差
\end{itemize}

\textbf{目标}：跟踪期望轨迹。

\subsection{导纳控制}

在导纳控制中，我们使用当前状态：
\begin{equation}
M\ddot{x}_d + B\dot{x} + Kx = f_{\text{ext}}
\end{equation}

这里：
\begin{itemize}
\item $x$: 当前实际位置（不是误差）
\item $\dot{x}$: 当前实际速度（不是误差）
\item 没有明确的期望位置$x_d$
\end{itemize}

\textbf{目标}：响应外力，呈现阻抗特性。

\section{总结}

\subsection{关键要点}

\begin{enumerate}
\item \textbf{导纳控制中使用当前状态}：
\begin{itemize}
\item $x$: 当前实际位置
\item $\dot{x}$: 当前实际速度
\item 不是误差项（不是$x_d - x$）
\end{itemize}

\item \textbf{物理意义}：
\begin{itemize}
\item 使机器人表现得像阻抗系统
\item 阻抗系统的力基于当前状态
\item 因此使用当前状态是合理的
\end{itemize}

\item \textbf{与轨迹跟踪的区别}：
\begin{itemize}
\item 轨迹跟踪：使用误差项，目标是跟踪期望轨迹
\item 导纳控制：使用当前状态，目标是响应外力
\end{itemize}
\end{enumerate}

\subsection{公式记忆}

\textbf{导纳控制公式}：
\begin{equation}
M\ddot{x}_d + B\dot{x} + Kx = f_{\text{ext}}
\end{equation}

\textbf{记忆方法}：
\begin{itemize}
\item $x$和$\dot{x}$是\textbf{当前实际状态}（测量值）
\item $\ddot{x}_d$是\textbf{期望加速度}（计算值）
\item 公式表示：虚拟力（基于当前状态）等于外力
\end{itemize}

\end{document}

